\section{Prescribed trajectories and their cost}\label{sec:paths}
\label{ac:part:18}

Fix a power of two \(R\ge2\), to be chosen sufficiently large. Let
\(\Omega=\Omega_1\times\cdots\times\Omega_b\) be an ordered product grid with
\(b\ge1\), where \(m_j=|\Omega_j|\) is a power of two in \([R,R^2]\), and put
\(s=|\Omega|=\prod_jm_j\). Each \(\Omega_j\) is identified with
\(\{0,1\}^{\log_2m_j}\). A sweep first permutes the lines parallel to coordinate 1,
then those parallel to coordinate 2, and so on, independently on all lines.
A line always belongs to its particular stage.

On a line of size \(m\), use
\begin{equation}
 L_m(z)=(1-z)U_m+zB_m,
 \label{ac:eq:2}
\end{equation}
where \(U_m\) is uniform on the permutations of that line and \(B_m\) is a
binary sweep on its \(\log_2m\) bits. Both laws send each individual slot to a
uniform slot. Both also have expected sign zero: for \(B_m\), toggling one of
its independent fair switches reverses the sign. At \(z=1\), the full grid
sweep is the binary sweep with consecutive bits grouped into these
coordinates.

Until the analytic argument, \(z\) is real, nonnegative, and small. For the
fixed \(R\), let
\[
 A_R=\max_m\max_{g\in S_m}m!B_m(g),\qquad
 z_{\rm cmp}(R)=\min\{1/2,1/(A_R+1)\},
\]
where \(m\) ranges over the allowed line sizes. There are finitely many such
sizes, so \(z_{\rm cmp}(R)>0\). For \(0\le z\le z_{\rm cmp}(R)\),
\begin{equation}
 \tfrac12U_m\ \le L_m(z)\ \le 2U_m
 \label{ac:eq:3}
\end{equation}
as measures.

The order of the coordinates determines a path from its endpoints. If a card
starts at \(u=(u_1,\ldots,u_b)\) and ends at \(v=(v_1,\ldots,v_b)\), then its
slot after stage \(j\) must be
\[
 (v_1,\ldots,v_j,u_{j+1},\ldots,u_b),\qquad 0\le j\le b.
\]
We may therefore prescribe a family \(\mathcal H\) of \(h\) such paths.
We require their slots to be distinct at every layer, including input and
output. These occupied slots will be called holes. For a stage line \(L\) of
size \(m\), let \(h_L\) count the prescribed paths using it, and let
\(E_L(\mathcal H)\) be the event that the line permutation realizes their
assignments from input to output. Layer disjointness makes these assignments
injective, and \eqref{ac:eq:3} makes \(E_L(\mathcal H)\) have positive
probability. The conditioning event is the intersection of these events over
independent lines. Thus the residual line laws are the original laws
\(L_m(z)\) conditioned on their prescribed assignments, and they remain
independent. Each acts as a random bijection between the unoccupied input and
output slots of its line.

Put \(M=s-h\). After the paths in \(\mathcal H\) are removed, the sweep gives
a random bijection between the \(M\) remaining input and output slots. For
\(\lambda\vdash M\), let \(K_\lambda^{\mathcal H}\) be its average in
\([\lambda]\), using any fixed identifications of the two slot sets as in
Section~\ref{sec:prelim}.

For an integer \(0\le u\le m\), define the line cost
\[
 \phi_m(u)=\log\frac{m^u}{(m)_u},\qquad
 (m)_u=m(m-1)\cdots(m-u+1),\qquad (m)_0=1,
\]
and set
\[
 C(\mathcal H)=\sum_{\text{stage lines }L}\phi_{|L|}(h_L).
\]
The summands in
\(\phi_m(u)=\sum_{r=0}^{u-1}\log(m/(m-r))\) are nonnegative and increasing.
For \(u\ge1\), their mean over the first \(u\) terms is at most their mean
over all \(m\) terms, which is at most one because \(m!\ge(m/e)^m\).
The case \(u=0\) is immediate. Hence \(0\le\phi_m(u)\le u\), and
\[
 0\le C(\mathcal H)\le bh,
\]
since the \(h_L\) sum to \(h\) at each stage.

For a nonnegative integer \(t\), a placement of \(t\) distinguishable cards is an injection from
\([t]=\{1,\ldots,t\}\) into the available slots; for \(t=0\) there is one
empty placement. Let \(x,y\) be placements in the remaining input and output
slots, and write \(p_{\mathcal H}(x,y)\) for their conditional transition
probability. If this probability is positive, the endpoint pairs determine
additional paths that, together with \(\mathcal H\), form a disjoint family
\(\mathcal H^+\) of size \(h+t\). Let \(v_L\) count the added paths on line
\(L\), so \(h_L+v_L\le |L|\). Figure~\ref{fig:trajectory} shows how a
prescribed path and an added path can share a line while occupying distinct
slots at every layer.

\begin{figure}[H]
\centering
\begin{tikzpicture}[x=1cm,y=1cm,>=Latex]
\foreach \x in {0,4.5}{
 \path[fill=gray!12,draw=gray!55,rounded corners=3pt]
   ({\x-.16},-.16) rectangle ({\x+.16},1.56);
}
\foreach \x/\name in {0/{input},4.5/{after coordinate 1},9/{output}}{
 \node[font=\small] at ({\x+.5},2.65) {\name};
 \foreach \y in {0,.7,1.4}
   \draw[gray!40] (\x,\y)--({\x+1},\y);
 \foreach \dx in {0,1}
   \draw[gray!40] ({\x+\dx},0)--({\x+\dx},1.4);
 \foreach \y in {0,.7,1.4}
   \foreach \dx in {0,1}
     \fill[gray!65] ({\x+\dx},\y) circle (1.2pt);
}
\node[font=\scriptsize,align=center] at (2.75,1.95)
  {coordinate 1\\vertical lines};
\node[font=\scriptsize,align=center] at (7.25,1.95)
  {coordinate 2\\horizontal lines};

\draw[very thick,blue!70!black,->,shorten >=2pt,shorten <=2pt]
  (0,1.4) to[out=-8,in=172] (4.5,.7);
\draw[very thick,blue!70!black,->,shorten >=2pt,shorten <=2pt]
  (4.5,.7) to[out=15,in=165] (10,.7);
\draw[thick,dashed,orange!85!black,->,shorten >=2pt,shorten <=2pt]
  (0,.7) to[out=-8,in=172] (4.5,0);
\draw[thick,dashed,orange!85!black,->,shorten >=2pt,shorten <=2pt]
  (4.5,0) to[out=-15,in=195] (10,0);

\foreach \p in {(0,1.4),(4.5,.7),(10,.7)}
 \fill[blue!70!black] \p circle (2.1pt);
\foreach \p in {(0,.7),(4.5,0),(10,0)}
 \node[rectangle,fill=orange!85!black,inner sep=1.7pt] at \p {};
\node[font=\small,text=orange!85!black] at (.5,-.75) {$(i,j)$};
\node[font=\small,text=orange!85!black] at (5,-.75) {$(i',j)$};
\node[font=\small,text=orange!85!black] at (9.5,-.75) {$(i',j')$};
\end{tikzpicture}

\caption{Schematic of selected slots in three layers of a grid with two coordinates.
The first coordinate is drawn vertically and the second horizontally.
The solid prescribed path and dashed added path share the highlighted
line at the first stage but use distinct slots at every layer. The added path from
\((i,j)\) to \((i',j')\) must pass through \((i',j)\). On the shared line this
configuration has \(h_L=1\) and \(v_L=1\), illustrating how a prescribed
assignment reduces the slots available to an added path.}
\label{fig:trajectory}
\end{figure}

On a line used by the added paths, with \(m=|L|\) and \(v=v_L>0\), the
conditional probability is
\[
 \frac{L_m(z)(E_L(\mathcal H^+))}
      {L_m(z)(E_L(\mathcal H))}
 \le \frac4{(m-h_L)_v}.
\]
Indeed the corresponding uniform ratio is exactly \((m-h_L)_v^{-1}\);
\eqref{ac:eq:3} changes its numerator and denominator by factors between
\(1/2\) and \(2\). The uniform ratio records the increment of the line cost:
\[
 \frac1{(m-h_L)_v}
 =m^{-v}\exp\{\phi_m(h_L+v)-\phi_m(h_L)\}.
\]
Unused lines contribute one, and at most \(tb\) lines are used. Each added
path contributes \(\prod_jm_j^{-1}=s^{-1}\). Multiplying the line bounds gives
\begin{equation}
 p_{\mathcal H}(x,y)
 \le s^{-t}\exp\{C(\mathcal H^+)-C(\mathcal H)+(\log4)tb\}.
 \label{ac:eq:4}
\end{equation}
The difference of potentials is the part of this estimate that permits
further paths to be revealed during the induction.

For a positive integer \(q\) to be fixed, define the unnormalized Schatten
moment
\[
 T_\lambda^{\mathcal H}=\|K_\lambda^{\mathcal H}\|_{2q}^{2q}.
\]

\begin{samepage}
\begin{theorem}[Conditional sweep moment estimate]
There\label{thm:conditional}\label{ac:main} are a power of two \(R\ge2\), a positive integer \(q\), and \(z_*>0\) such
that every grid and every family \(\mathcal H\) of prescribed disjoint
trajectories defined above satisfy, for every \(z\in[0,z_*]\) and every
\(\lambda\vdash s-h\),
\begin{equation}
 \log T_\lambda^{\mathcal H}\ \le\
 -c(s)F(\lambda)+e(s)h\log s-C(\mathcal H).
 \label{ac:eq:15}
\end{equation}
Here a zero moment gives \(-\infty\). Use parameters
\[
 a=\frac1{100},\quad c_0=e_0=\frac a{100},\qquad
 c(s)=c_0+\frac1{\sqrt{\log s}},\qquad
 e(s)=e_0-\frac1{\sqrt{\log s}}.
\]
The empty diagram has dimension one. If \(h=s\), the remaining bijection is
the unique empty bijection and its moment is one. The assertion also includes
the one-dimensional trivial representation for every \(h\).
\end{theorem}
\end{samepage}

For \(\mathcal H=\varnothing\), the estimate is
\(T_\lambda^\varnothing\le D_\lambda^{-c(s)}\); Section~\ref{sec:analytic}
uses this estimate for small real parameters in the analytic transfer. For general
\(\mathcal H\), the term \(-C(\mathcal H)\) is retained because
\eqref{ac:eq:4} expresses the cost of prescribing additional paths as a difference of
potentials. Section~\ref{sec:induction} uses that difference to recover the
original potential after summing over the added placements.
